% Chapter 1

\chapter{INTRODUCTION}\label{cha:intro}

Unmanned aerial vehicles (UAVs) are a practical option for last-mile delivery where
ground transport is slow or impossible: medical supplies for remote clinics, emergency
drug drops after a disaster, and diagnostic samples that should not spend an hour
shaking in the back of a van. In each case the UAV flies from a dispatch hub to a
delivery point through airspace that has buildings, terrain, restricted zones and wind
in it. A good route has to balance three things that pull against each other: how far
the drone flies, how much battery it uses, and how much of the flight is spent in risky
air.

Multi-objective path planning was built for this kind of conflict. A multi-objective
planner keeps distance, energy and risk separate during the search and returns a
Pareto-optimal front, which is a set of routes where none can be improved on one
objective without getting worse on another~\cite{r11}. The front tells the operator
which trade-offs exist. It does not say which route to fly.

The usual shortcut is to fold the three objectives into one weighted sum and run an
ordinary planner such as A* on the resulting scalar cost. The operator picks the
weights before each flight to express what the mission cares about. For one fixed
mission type this works. It stops working once the payload changes, because emergency
medicine needs a very different weight profile from a fragile blood sample or a box of
relief supplies~\cite{r1}. Re-tuning the weights and running the search again between
missions is slow, needs someone who understands the planner, and is hard to justify in
the middle of an emergency.

This project proposes a Payload-Aware Two-Tier UAV Route Planning system that removes
that step. Tier~1 builds the Pareto front for a delivery region once, with an exact
multi-objective graph search over a three-dimensional airspace lattice whose edge costs
depend on terrain, buildings, restricted zones and the direction of the wind. The front
is stored and reused until the graph changes, the front expires or the weather drifts.
Tier~2 runs per mission. It reads four payload tags entered by the operator, turns them
into objective weights with a fixed formula, scores every route on the stored front and
returns the best one in well under a millisecond. The chosen route is exported as an
ArduPilot mission and flown in a Gazebo and ArduPilot software-in-the-loop (SITL)
simulation.

This report covers Phase~I: the software system, its benchmarks against
single-objective A* and MOPSO, and the first simulated flights. Phase~II moves the same
software onto a physical quadcopter.

\section{MOTIVATION}

Three delivery scenarios that come up again and again in medical and emergency drone
work motivate this project, and their routing needs do not agree with each other.

Fragile medical samples such as tissue biopsies and blood specimens must arrive intact,
and an operator will want them kept out of rough air when a smoother corridor exists.
The priority is low flight risk: little time in turbulent or crosswind-exposed corridors
and few sudden altitude changes. A somewhat longer but much calmer route is better.

Emergency medicine such as insulin or antivenom has a strict time window. The priority
is the shortest route, so that transit time is as low as possible, and risk and energy
come second. If the delivery window is fifteen minutes, the shortest path wins.

Disaster-relief supplies such as water purification tablets and ration packs are
heavier than medical payloads, and they are neither fragile nor urgent. Here the
priority is low battery drain, so that one charge covers as many drops as possible.

A real delivery service sees all three, sometimes within the same hour. A planner with
fixed weights treats them identically, and one that lets the operator change the weights
pushes an optimisation task onto someone dispatching a drone in a hurry. We could not
find a UAV delivery planner that
switches between these routing strategies from a description of the payload, without
an expert in the loop and without running the search again. That is what motivated the
two-tier design.

A second motivation is wind. None of the planners we reviewed resolve wind direction
per edge, yet for a quadcopter cruising at 36\,km/h a 25\,km/h wind is not a detail:
flying into it costs far more energy than flying with it, and crosswind makes the ride
rough.
Section~\ref{sec:res-wind} shows that at a fixed wind speed, the direction alone changes
the energy of the cheapest available route by 27.5\%. A planner blind to direction gets this
wrong on any windy day, and so does any payload-aware choice built on top of it.

\section{BACKGROUND}\label{sec:background}

\subsection{Building a front versus choosing from it}

Pareto fronts for routing are built either by exact graph search, such as
NAMOA*~\cite{r17}, which keeps every non-dominated partial path at each node, or by
population methods such as NSGA-II~\cite{r8} and MOPSO~\cite{r6}, which are faster on
large spaces but stochastic. Building the front is well studied. Picking one route from
it gets far less attention: the usual options are a crowding-distance rule~\cite{r8,r3},
an entropy criterion followed by a hand-set weighted sum~\cite{r1}, or a person. None of
them looks at what the drone is carrying. Payload does appear in the literature as
continuous mass inside the energy model~\cite{r10}, which changes the costs, not the
choice between trade-offs. What is missing is any use of payload metadata to
\emph{select} a point on a precomputed front.

\subsection{Evidence from real deployments}

Real deployments show where payload matters, and also how thin the public evidence
still is. Swiss Post flew laboratory samples between hospitals on Matternet quadrotors
until two incidents in 2019 led it to suspend the service. In January a short circuit
cut power to the GPS and the drone made an emergency landing on Lake Z\"{u}rich. In May
a parachute tether snagged on the airframe and the 22-pound drone fell into woods near a
group of kindergarten children~\cite{r12}. Both were equipment failures, so they say
nothing about route choice. They do show that laboratory samples are routinely flown
over water and populated ground, where the choice of corridor has consequences.

Laboratory studies of medicines and blood samples flown by drone found no damage on
ordinary flights~\cite{r13,r14,r15}, but none of them varied the route
(Section~\ref{sec:lit-papers}), so in this work the fragility tag is a precaution the
operator sets, not a measured degradation effect.

Zipline's blood-delivery network in Rwanda serves hospitals up to about 80\,km from its
distribution centre, and each aircraft follows a predetermined flight plan~\cite{r16}.
Its engineers treat payload weight as a limit on range. That is the energy side of the
problem Tier~2 addresses: a heavier load shortens how far the drone can go, so the best
route depends on what is on board.

\subsection{Regulation}

In India, small civil drones fly under the Drone Rules, 2021~\cite{r21}, which treat
120\,m above ground as the ceiling for flight without special permission. Ground risk is
commonly ranked with the SORA framework~\cite{r22}. Our airspace model uses both.

\section{PROBLEM DEFINITION}\label{sec:problem}

Existing UAV route planners either apply a fixed geometric selection rule over the
Pareto front that ignores the payload, or fold the objectives into one hand-tuned
weighted sum before the search begins, which fixes the trade-off regardless of what is
being carried. When the payload changes, an expert has to re-tune the weights and run the search again,
because no explicit mapping exists from a payload description to the trade-off it
implies. In the field, that expert is not available.

\subsection{Formal statement}

We are given:
\begin{enumerate}
  \item an airspace graph $G = (V, E)$ in which every edge carries a three-objective cost
  vector $\mathbf{c}(e) = [d(e),\, \varepsilon(e),\, \rho(e)]$ for flight distance,
  estimated energy and flight risk;
  \item a Pareto front $\mathcal{P}$ of non-dominated routes computed in advance over
  $G$; and
  \item a payload metadata tuple $(F, W, U, P_\mathrm{lvl})$ holding fragility class
  $F \in \{1,\ldots,5\}$, payload weight class $W \in \{1,2,3\}$, delivery urgency
  $U \in \{1,2,3\}$ and priority level $P_\mathrm{lvl} \in \{1,\ldots,5\}$.
\end{enumerate}
The problem is to select the route $p^* \in \mathcal{P}$ that best meets the
requirements implied by the payload metadata, without expert re-tuning and without
running the multi-objective search again.

The inputs are values an operator can assign from domain knowledge (``this vial is
fragile, $F = 5$; the delivery is not urgent, $U = 1$''). Selection must be
deterministic and reproducible, and it must finish within one second of the tags being
entered. Three further conditions follow from the gaps identified in
Chapter~\ref{cha:lit}. The selected route must depend on $(F, W, U, P_\mathrm{lvl})$ and
not on the geometry of the front alone. The front must be reused across missions
between refreshes instead of being rebuilt each time. The map from
$(F, W, U, P_\mathrm{lvl})$ to $p^*$ must be an explicit formula, not a set of rules
written by an expert and not a lookup table.

\subsection{Weather dependence}

The edge costs depend on the weather. Writing $\omega$ for a weather snapshot (wind
speed, wind direction, gust, rain and visibility), the cost of an edge is
$\mathbf{c}(e;\omega) = [d(e),\, \varepsilon(e;\omega),\, \rho(e;\omega)]$ and the
front becomes $\mathcal{P}(\omega)$. Distance stays the same, but energy and risk
change with how the wind meets each edge. The selection problem itself does not change.
What changes is that a stored front is only valid while $\omega$ stays close to the
snapshot it was built for. That is why the front is refreshed on a schedule and on
weather drift. When $\omega$ crosses the go/no-go limits, no route is selected at all
and the mission is held.

\subsection{Inputs, outputs and success criteria}

Table~\ref{tab:ps-io} lists what the system takes in, what it returns and how each part
is judged.

\begin{table}[!htbp]
\centering
\caption{Inputs, outputs, constraints and success criteria}\label{tab:ps-io}
\small
\begin{tabular}{|P{3.3cm}|P{11.2cm}|}
\hline
\textbf{Item} & \textbf{Definition}\\
\hline
Region input & Centre point and box size; terrain (DEM), OpenStreetMap buildings, restricted zones, water and land use for that box\\
\hline
Weather input & Snapshot $\omega$ from Open-Meteo: wind at 80\,m, gust, rain, visibility, WMO weather code\\
\hline
Mission input & Depot, destination and payload tuple $(F, W, U, P_\mathrm{lvl})$\\
\hline
Tier~1 output & Stored Pareto front $\mathcal{P}(\omega)$ with graph hash, weather snapshot and expiry time\\
\hline
Tier~2 output & Selected route $p^*$ with waypoints (latitude, longitude, altitude), weights, ranks and point type, or a HOLD advisory\\
\hline
Mission output & QGC WPL~110 \texttt{.waypoints} and QGroundControl \texttt{.plan} files\\
\hline
Hard constraints & Cruise at least 60\,m above ground and 15\,m above roofs; at most 120\,m above ground; no entry into restricted zones; route energy within 80\% of battery capacity\\
\hline
Tier~1 criteria & Hypervolume and IGD at least as good as every baseline on each scenario; no dominated or duplicate routes in the front\\
\hline
Tier~2 criteria & Selection under 1\,s; zero monotonicity violations over all 225 tag combinations; non-extreme picks for mixed payloads\\
\hline
Simulation criteria & Cross-track error, clearance margin and landing offset of the flown route in SITL and Gazebo\\
\hline
\end{tabular}
\end{table}

\subsection{Scope}

Phase~I, covered by this report, delivers the full software system and its simulation.
Phase~II moves the system onto a Pixhawk-equipped
F450 quadcopter with a Raspberry Pi~5 companion computer and validates it with outdoor
flights. The Tier~1 front is treated as valid for a bounded window. It is rebuilt when
the window expires or when the weather drifts past set thresholds, and every mission
inside the window is served by Tier~2 alone. The project does not cover multi-UAV
coordination, avoidance of moving obstacles during flight, or real-time replanning.

\section{ORGANIZATION OF THE REPORT}

The report has six chapters, followed by the references and three appendices. Chapter~\ref{cha:intro} sets out the motivation, background and
problem. Chapter~\ref{cha:lit} reviews the closest papers, the gaps in them and our
research objectives. Chapter~\ref{cha:method} describes the architecture and its seven
modules with the algorithms behind them. Chapter~\ref{cha:results} covers the datasets,
the hardware and software, the experiments and what the results show. Chapter~\ref{cha:social} discusses the social impact and sustainability
of the work. Chapter~\ref{cha:conclusion} concludes and lists the work planned for the
rest of Phase~I and for Phase~II. The references follow, and the appendices give the
main configuration parameters, the automated test suites and the code organisation.
